\subsection{Literature Survey}
\label{sec:literature}

\paragraph{Recurrent networks on raw light curves.}
\citet{charnock2017} were among the first to show that a recurrent neural
network fed raw, un-engineered light-curve photometry directly could
classify supernova types competitively with feature-engineering pipelines,
on simulated Dark Energy Survey-like data -- establishing that a sequence
model need not first be handed hand-built summary statistics (rise time,
colour at peak, and so on) to find useful structure in a light curve.
\citet{muthukrishna2019rapid} extended this to \emph{early}, time-resolved
classification with RAPID, a Gated Recurrent Unit (GRU) architecture that
emits a classification at every observed timestep rather than only at the
end of a light curve, across 12 transient classes; this per-timestep
prediction convention, and the resulting practice of reporting accuracy as
a function of time since first detection rather than as one end-of-light-
curve number, is used throughout the present work (Section~\ref{sec:analysis}).
\citet{moller2020supernnova} (SuperNNova) placed a Bayesian recurrent
architecture in this same family, adding calibrated uncertainty estimates
and a systematic comparison against feature-based random-forest baselines,
run at ZTF and LSST alert-stream scale.

\paragraph{Data challenges.}
\citet{kessler2019plasticc} designed and described PLAsTiCC, the Photometric
LSST Astronomical Time-series Classification Challenge: a large, publicly
released simulated dataset spanning many transient and variable-source
classes, intended to let classification methods be compared on a common,
labelled corpus before real LSST data existed. ELAsTiCC2 \citep{elasticc2},
the LSST Dark Energy Science Collaboration's Extended LSST Astronomical
Time-series Classification Challenge and this project's own Tier A
simulated corpus (Section~\ref{sec:data}), supersedes PLAsTiCC with an
updated simulation
pipeline, revised class definitions, and -- the property this project
depends on most directly -- deliberately \emph{non}-representative class
proportions relative to the real sky rate, so that a classifier trained on
it must be evaluated with that mismatch in mind rather than assumed to
inherit a realistic prior (Section~\ref{sec:data}).

\paragraph{Current state of the art.}
\citet{cabreravives2024atat} report ATAT, a pair of transformer encoders
(one over the light curve, one over tabular host/metadata features) fused
before a classification head, reaching a macro-F1 of $82.9 \pm 0.4$ across
20 ELAsTiCC classes -- the strongest published number on this class of
problem to date, and the anchor this project's own, considerably smaller
3-class macro-F1 figures should be read against, not compared to directly
(Section~\ref{sec:results}). \citet{shah2025oracle} (ORACLE) instead build
a GRU-based hierarchical classifier trained with a hierarchical
cross-entropy loss and host-galaxy context, explicitly targeting the same
graceful-degradation property (a correct ``supernova-like'' call even when
the fine-grained subtype call is wrong) that this project's own hierarchical
loss experiment tests and does not, in the end, adopt
(Section~\ref{sec:methodology}). A causal-transformer architecture for this
same task is described in \citet{atcat2025}, reporting 71.8\% light-curve-only
and 89.8\% light-curve-plus-metadata F1 on their own evaluation. Fine-grained
supernova subtyping is consistently reported, across this literature, as a
harder discrimination axis than the coarser transient-vs-non-transient call
-- consistent with this project's own recorded expectation that the Ib/c
class in particular should be, and is, the hardest of the three target
classes to recover cleanly.

\paragraph{Anomaly detection over alert streams.}
\citet{malanchev2021} ran an unsupervised anomaly detector (isolation-forest-
based) over ZTF's third public data release and reported that of the
highest-ranked, most anomalous objects a human inspected, 68\% were
artifacts of imperfect image subtraction rather than genuine unusual
transients -- the published anchor this project's own Gate D (deferred, not
passed -- Section~\ref{sec:results}) was designed to be checked against.
\citet{villar2021vraenn} (VRAENN) instead trained a variational recurrent
autoencoder for the same live-anomaly-detection task, an architecture whose
reconstruction-error principle is close in spirit to the decoder component
of this project's own three-part anomaly score
(Section~\ref{sec:methodology}). \texttt{coniferest} \citep{coniferest}, the
isolation-forest and active-learning package used here for one of the three
anomaly components, is developed and maintained by the SNAD (SNAD Anomaly
Detection) team; the production Fink broker's own anomaly-ranking module,
described in \citet{moller2021fink}, is built on the same package.

\paragraph{Real data and brokers.}
The Zwicky Transient Facility Bright Transient Survey (BTS,
\citealp{fremling2020bts,perley2020bts}) is the source of this project's
real, spectroscopically labelled fine-tuning data (Tier B). BTSbot
\citep{rehemtulla2024btsbot} is a real-bogus and bright-transient-vs-reject
triage classifier trained on the same BTS alert stream; this project reuses
its public metadata release as the raw Tier B corpus (Section~\ref{sec:data})
but not its own triage labels, which answer a different question (``is this
a real, bright transient at all'') from the one this project trains on
(``which of Ia/Ib-c/II is it''). ALeRCE \citep{forster2021alerce} and Fink
\citep{moller2021fink} are both production alert brokers that ingest ZTF
(and, for Fink, LSST) alert streams and run their own classification and
anomaly-ranking pipelines in real time; Fink is the broker this project's
own serving layer (Section~\ref{sec:methodology}) actually connects to for
live Tier C data.

\paragraph{Transfer learning.}
\citet{guptamuthukrishna2025} report that pretraining a transient classifier
on simulated data and fine-tuning on a comparatively small real labelled
ZTF set reduces the amount of real labelled data needed by about 75\%
relative to training on real data alone, and that a model pretrained on
ZTF reaches 94\% of a from-scratch baseline's performance on simulated LSST
data using only 30\% as much LSST training data -- the same sim-to-real
transfer strategy this project's own Tier A~$\to$~Tier B pipeline follows
(in the opposite direction, LSST-like simulation to real ZTF), and the same
comparison (zero-shot Tier-A-only performance versus fine-tuned performance
on Tier B) this project measures directly in Section~\ref{sec:results}.

\paragraph{Loss design.}
\citet{villar2023whxe} propose weighted hierarchical cross-entropy (WHXE): a
loss that predicts at every level of a class taxonomy tree, not only at the
leaf, and weights a mistake by how far apart the true and predicted classes
sit in that tree -- confusing two supernova subtypes should cost less than
confusing a supernova with a non-supernova. This project implements WHXE
essentially as published, encounters and resolves a real ambiguity in the
paper's own depth-indexing convention, and reports a controlled, one-run
comparison against a flat classification head that does not favour WHXE on
this project's own single-branch dataset (Section~\ref{sec:methodology});
this is stated as a property of the data available at the time of the
comparison, not as a criticism of the loss itself.

\paragraph{Surveys.}
\citet{bellm2019ztf} describe ZTF's system design, cadence, and early
performance; \citet{ivezic2019lsst} describe LSST's science drivers,
reference design, and anticipated data products, including the alert-volume
figures cited in Section~\ref{sec:introduction}.

\paragraph{Taxonomy.}
\citet{filippenko1997} is the standard review of the spectroscopic
classification scheme this entire field, and this project's own taxonomy
(Section~\ref{sec:data}), is built on.

\paragraph{Where this work sits.}
None of the architectural components used here are novel: a causal GRU
light-curve encoder, a reconstruction-error/Mahalanobis/isolation-forest
anomaly ensemble, and a Kafka-based serving layer are all established
techniques, individually published elsewhere in the works cited above. What
this report adds is not a new architecture but a measured, audited account
of what happens when a comparatively standard pipeline of this kind is
actually built end to end on a strict timeline: a quantified sim-to-real
transfer gap and transfer gain (Section~\ref{sec:results}); a documented
account of where the train-time and serve-time data distributions silently
diverge and what that costs (Section~\ref{sec:analysis}); a measured,
non-obvious failure mode in the anomaly ensemble under domain shift
(Section~\ref{sec:results}); and end-to-end serving and live-ingestion
results, run against real ZTF and LSST alert traffic, that are not published
elsewhere for a pipeline of this specific shape. Understating this
contribution is the safer error to make, and is made deliberately.
